\section*{ЗАКЛЮЧЕНИЕ}

Работа посвящена поддержанию актуального соответствия
\textit{IP-префикс~$\rightarrow$~ASN} по данным BGP для продуктов BIFIT
MITIGATOR и BIFIT COLLECTOR. Главный результат -- переход от
периодической полной перестройки таблицы к потоковому обновлению
in-memory состояния в отдельном микросервисе.

\textbf{Потоковое обновление in-memory базы.} Реализована операционная
модель «начальный снимок + поток событий» от узла GoBGP: обновления
применяются к одной общей копии состояния по одному событию, с
автоматической пересинхронизацией после разрывов. На материале работ по
динамике BGP~\cite{labovitz1997instability,labovitz2000convergence,
griffin2002stablepaths} это согласуется с непрерывным характером
маршрутного потока и несовместимостью с одной лишь периодической полной
перестройкой.

\textbf{Вынесение функции в микросервис.} Логика построения и обновления
таблицы перенесена из состава BIFIT COLLECTOR в автономный процесс с
единым gRPC API для нескольких потребителей. Это устраняет дублирование
между продуктами и фиксирует границу ответственности: сопоставление по
данным BGP находится в одной подсистеме, независимо от выбранного
варианта хранилища
(\texttt{ranger} или шардированный ART).

\textbf{Структура хранения и корректность снимка.} После обзора
критериев для префиксных структур в главе~\ref{sec:storage-choice} в качестве
технической основы потока выбран и реализован шардированный
Adaptive~Radix~Tree~\cite{leis2013art} с операциями \texttt{Replace},
\texttt{Upsert}, \texttt{Delete}, \texttt{Lookup} и \texttt{Dump}. Для
внешних клиентов доступны \texttt{GetASNInfo},
\texttt{GetPrefixesASTable} и \texttt{SubscribePrefixASUpdates}. Полный
снимок строится без перекрывающихся префиксов в выдаче.

\textbf{Эксперимент и компромиссы.} Сравнение на одном стенде двух
реализаций одного \texttt{prefixStore}-контракта показало: при
$n = 10^5$ шардированный ART быстрее \texttt{ranger} по
\texttt{Upsert}, а \texttt{Lookup} с фильтрацией по активным длинам
префиксов сопоставим у обеих реализаций. Полный
\texttt{Dump} дороже у ART, поскольку внешняя выдача строится как
непересекающаяся таблица. Отдельный прогон на $3\cdot10^6$ смешанных
IPv4/IPv6 записей усиливает этот вывод: ART выполняет \texttt{Replace}
примерно за $693{,}98$~мс против $30{,}84$~с у \texttt{ranger},
\texttt{Lookup} -- за $176{,}7$~нс против $187{,}8$~нс, а
\texttt{Upsert} -- за $255{,}5$~нс против $7{,}99$~мкс; при этом
\texttt{Dump} занимает $6{,}16$~с против $4{,}25$~с у \texttt{ranger}.

Сопоставление с задачами работы: главы~\ref{sec:bgp-domain}--\ref{sec:bgp-dynamics}
дают теоретическую базу; глава~\ref{sec:storage-choice} обосновывает
выбор хранилища; глава~\ref{sec:problem} формулирует требования и
гипотезу; главы~\ref{sec:architecture}--\ref{sec:art-storage} описывают
архитектуру микросервиса, алгоритмы снимка и потока, реализацию
хранилища; главы~\ref{sec:methodology}--\ref{sec:experiment-results}
содержат методику и численные оценки.

Рекомендации: как единый источник пары «префикс~$\rightarrow$~ASN» в
режиме реального времени сервис с шардированным ART предпочтителен; если
доминирует частая полная выгрузка большой таблицы без потока, нужен
\texttt{ranger} или кэшированный/потоковый \texttt{Dump}.

Перспективы: дальнейшая оптимизация LPM поверх ART, кэшированный
упорядоченный снимок для \texttt{Dump}, интеграция с
BMP~\cite{rfc7854} и коллекторами
\cite{ripe-ris,routeviews,ripe-beacons}, проверки на данных, близких к
промышленным.

%==============================================================
